% ===================================================================
%  HOTO Na 3 — Ƙuruciya da Samartaka.
%
%  Traduction, faite en 2026, en haoussa d'aujourd'hui, sur l'ido de
%  texto/io. Regles de la colonne : voir 10-hoto-01.tex. Deux scenes,
%  deux numerotations : t03-c1-* le bapteme au village, t03-c2-* la
%  fete publique.
%
%  LE « julio » DU BLOC t03-c2-01-1. Le fac-simile ido compose
%  « (julio od \VUgras{agosto}) » en laissant le seul juillet hors du
%  gras, alors que les deux mois qui l'encadrent y sont.
%  gravuri/korekti.json le repare pour chaque colonne, et le triple
%  haoussa — « (Yuli ko » — a ete pose AVANT que ce tableau soit
%  ecrit : la correction attend le texte, non l'inverse.
%
%  L'ANNEE HAOUSSA A CINQ SAISONS, ET LA PLANCHE N'EN CONNAIT QUE
%  QUATRE. Le Sahel se partage en rani la saison seche et chaude,
%  bazara la saison qui la suit et prepare la pluie, damina la saison
%  des pluies, kaka la moisson, hunturu le froid de l'harmattan. Le
%  printemps de ce tableau tombe juste : bazara est bien mars-avril-
%  mai, la saison des bourgeons. L'ete de la seconde scene est rani,
%  comme au tableau 2. Mais l'automne et l'hiver du tableau 4 auront
%  a choisir entre kaka et hunturu, et ce sera le seul endroit du
%  livret ou une saison europeenne n'ait pas de case ou tomber.
%
%  LES MOIS SONT ANGLAIS ET LES JOURS SONT ARABES, ET LA LANGUE
%  N'HESITE PAS UNE FOIS. Janairu, Faburairu, Maris, Afirilu, Mayu,
%  Yuni, Yuli, Agusta, Satumba : le calendrier gregorien est venu par
%  Lagos, et « Afirilu » montre au passage la loi du p — l'alphabet
%  boko n'en a pas, et April devient Afirilu. Les jours, eux, sont
%  arabes tous les sept (voir le tableau 2), parce que la semaine est
%  venue avec l'islam. Et le calendrier VECU ne porte ni les uns ni
%  les autres : le mois du jeune s'appelle « Watan Azumi », le mois
%  du jeune ; celui de la fete, « Watan Sallah ». Ce qu'on emprunte
%  garde son nom etranger ; ce qu'on vit se nomme par ce qui s'y
%  fait.
%
%  LA MOSQUEE EST ARABE ET L'EGLISE EST ANGLAISE, ET C'EST LA MEME
%  FRONTIERE QUE CHEZ LE PERSAN, TRACEE AVEC D'AUTRES LANGUES. Le
%  persan disait غسل تعمید pour le bapteme, کشیش pour le pretre,
%  کلیسا pour l'eglise — des mots arabes ou syriaques qu'il ne
%  confond jamais avec مسجد et ملا. Le haoussa fait exactement le
%  meme partage : masallaci la mosquee, liman l'imam, ladan le
%  muezzin, sallah la priere, tous arabes ; coci l'eglise, baftisma
%  le bapteme, fada le cure, tous anglais. Deux religions, deux
%  lexiques, et pas un mot qui passe de l'un a l'autre. Une langue
%  qui a vecu mille ans a cote d'une autre religion sait la nommer
%  sans l'assimiler — le persan le montrait au meme tableau.
%
%  LE PAPILLON (35) SE DIT « malam buɗe littafi », C'EST-A-DIRE
%  « maitre, ouvre le livre ». C'est le plus beau nom du releve, et il
%  tient dans un imperatif : les deux ailes qui battent sont les deux
%  pages d'un livre qu'on ouvre et referme, et l'enfant qui le voit
%  passer appelle le malam du tableau 1. Aucune autre colonne ne
%  nomme un insecte par une phrase.
%
%  ET LA GIRAFE (53) EST LE « raƙumin dawa », LE CHAMEAU DE LA
%  BROUSSE. La ou le haoussa connait la bete, il la nomme par celle
%  qu'il connait mieux ; la ou il ne la connait pas, il emprunte et
%  ne cache rien : « beyar » l'ours au tableau 2, « damisa » ici pour
%  le tigre (51) — mais damisa est le leopard, et le haoussa n'a pas
%  de mot pour un fauve d'Asie. Le lion (52), lui, est « zaki », mot
%  d'ici, comme le شیر persan et pour la meme raison : la bete etait
%  la.
%
%  DEUX ARBRES QUI NE POUSSENT PAS AU SAHEL, DEUX FRUITS QUI N'Y
%  MURISSENT PAS, ET DEUX CONDUITES OPPOSEES. Le peuplier (3) et le
%  platane (17) n'ont pas de nom : la colonne les decrit par leur
%  port — « sirarun bishiyoyi », les arbres minces, « dogon icen
%  inuwa », le grand arbre d'ombre. La peche et l'abricot (18, 19),
%  eux, ont un nom, emprunte et refait a la bouche haoussa : fis,
%  abirkot. LE FRUIT ARRIVE DANS UNE CAISSE, L'ARBRE N'ARRIVE JAMAIS
%  — et la langue nomme ce qu'elle a tenu.
%
%  QUATRE FLEURS, QUATRE CHEMINS : l'oeillet (34) est « furen
%  kanumfari », la fleur de girofle, exactement comme le francais
%  l'appelait giroflee et l'arabe قرنفل ; la tulipe (36) est
%  « tulif », emprunt refait au p ; le muguet (37) est « furannin
%  Mayu », les fleurs de mai, ce que l'ido dit lui-meme ; la rose
%  (39) est « wardi », de l'arabe. Une fleur qu'on sent, une qu'on
%  achete, une qu'on cueille, une qu'on ecrit.
%
%  LE BAL PUBLIC (42) EST LE PREMIER REFUS DE LA COLONNE, ET C'EST LE
%  MEME REFUS QUE LE JAVANAIS A FAIT SUR LA MEME PLANCHE. Le javanais
%  avait « joged » et ne l'a pas ecrit ; le haoussa a « bori », qui
%  est la danse de possession, et il a aussi « asauwara » et les
%  rawa de metier. La planche grave un bal de sous-prefecture avec
%  estrade et lampions : la colonne ecrit « wurin rawa », le lieu ou
%  l'on danse, et rien de plus. On modernise la langue, jamais les
%  choses — et c'est la deuxieme fois que ce tableau-ci le demande.
%
%  UNE FAUTE D'ORDRE, LA PREMIERE DE LA COLONNE, ET C'EST EXACTEMENT
%  CELLE QUE LE JAVANAIS AVAIT FAITE — MEME TABLEAU, MEME BLOC, MEME
%  PHRASE. Au bloc t03-c2-07-1 le feu d'artifice (20) avait ete pose
%  avant les roseaux (19), alors que l'ido dit « preparas inter la
%  kani la pirotekno » : le lieu d'abord, l'objet ensuite. Les deux
%  colonnes qui se sont trompees la sont a tete initiale et posent
%  naturellement l'objet avant son lieu ; rien ne les y obligeait —
%  « cikin kyauro » s'insere tres bien avant le complement d'objet —,
%  c'etait un simple choix de tour. DEUX LANGUES SANS AUCUN RAPPORT
%  QUI GLISSENT AU MEME MOT : la faute n'est pas dans la langue, elle
%  est dans la PHRASE de l'ido, et seule la machine peut la voir.
%
%  ENFIN L'ONCLE (56) : le haoussa compte ses oncles comme le persan
%  comptait les siens. baffa est le frere du pere, kawu le frere de
%  la mere, goggo la soeur du pere, inna la soeur de la mere. La
%  planche ne dit pas de quel cote vient l'oncle ; la colonne a pris
%  baffa, et la tante (58) devient goggo pour rester du meme cote.
%  C'est le choix exact qu'avait fait la colonne persane avec عمو, et
%  pour la meme raison : une langue qui distingue oblige a savoir ce
%  que le texte d'origine n'a pas eu besoin de dire.
% ===================================================================

%%K t03-apar-1 apar
\VUsaut{3.0mm}
\VUtitre{15.0pt}{60}{HOTO Na 3}
\VUsaut{1.6mm}
\VUtitre{11.4pt}{0}{Ƙuruciya da Samartaka.}
\VUsaut{3.4mm}

%%K t03-apar-2 apar
(An haɗa hoto na 3 da na 4 ta yadda suke nuna, cikin \VUgras{fage
huɗu na iyali}, muhimman ra’ayoyi game da \VUgras{yanayi}, da
\VUgras{shekaru}, da \VUgras{iyali}, da \VUgras{tufafi} da
\VUgras{hali}.)

\VUsaut{4.0mm}
%%K t03-tit-1 sub
\VUcentre{12.0pt}{0}{\textit{Fage na Farko.}}
\VUsaut{3.0mm}

%%K t03-tit-2 sub
\VUpk{10.2pt}{40}{Bikin baftisma a ƙauye.}
\VUsaut{2.4mm}

%%K t03-tit-3 sub
\VUpk{10.2pt}{40}{Bazara.}
\VUsaut{3.0mm}

%%K t03-c1-01-1 p
1. --- Muna cikin \VUgras{bazara,} a watannin \VUgras{Maris,}
\VUgras{Afirilu} ko \VUgras{Mayu,} a ranakun \VUgras{mako,}
\VUgras{Litinin,} \VUgras{Talata,} ko \VUgras{Laraba.} \VUgras{Ƙarfe
goma} na safe ne.

%%K t03-c1-02-1 p
2. --- \VUgras{Yanayi} yana da kyau, \VUgras{iska} tana da tsabta.
Rana tana haskakawa. Waɗansu \VUgras{ƙananan gajimare}
\textsuperscript{(2)} suna hawa a cikin shuɗiyar \VUgras{sama}
\textsuperscript{(1)}.

%%K t03-c1-03-1 p
3. --- \VUgras{Bishiyoyi} kusan ba su da \VUgras{ganye.}
\VUgras{Sirarun bishiyoyi} \textsuperscript{(3)} da \VUgras{dogon icen
inuwa} \textsuperscript{(17)} har yanzu ba su da komai sai toho.

%%K t03-c1-04-1 p
4. --- \VUgras{Alallaka} \textsuperscript{(4)} suna dawowa suna
shawagi cikin ƙwarƙwasa mai farin ciki kewaye da \VUgras{tsinin}
\textsuperscript{(7)} \VUgras{hasumiyar ƙararrawa} \textsuperscript{(5)}
wadda take kambin \VUgras{cocin} \textsuperscript{(6)} ƙauyen.

%%K t03-tit-4 sub
\VUsaut{3.4mm}
\VUpk{10.2pt}{40}{Coci.}
\VUsaut{3.0mm}

%%K t03-c1-05-1 p
5. --- Wannan coci mai sauƙi ne ƙwarai, tare da \VUgras{babbar
ƙofarsa} \textsuperscript{(13)} da \VUgras{babban agogonsa}
\textsuperscript{(8)}. \VUgras{Ƙaramar allura} \textsuperscript{(9)}
tana kan lamba X, tana nuna \VUgras{sa’o’i.} \VUgras{Babbar}
\textsuperscript{(10)} tana kan XII (tsakar rana) ; tana nuna
\VUgras{mintuna.}

%%K t03-c1-06-1 p
6. --- A cikin hasumiya, \VUgras{ƙararrawa} \textsuperscript{(11)} suna
kaɗawa cikin farin ciki. A saman rufin akwai ƙaramin
\VUgras{gicciyen} \textsuperscript{(12)} dutse a tsaye.

%%K t03-c1-07-1 p
7. --- Coci ba \VUgras{babbar ƙofa} \textsuperscript{(13)} kaɗai yake
da ita ba, yana da ƙaramar ƙofa a gefe. Ta wannan ƙaramar ƙofa ne
malam \VUgras{fada} \textsuperscript{(15)}, sanye da \VUgras{doguwar
rigarsa,} yake fitowa daga \VUgras{ɗakin tufafi}
\textsuperscript{(14)} ya nufi \VUgras{gidan fada}
\textsuperscript{(19)}.

%%K t03-c1-08-1 p
8. --- Kusa da shi, ga \VUgras{mai sayar da madara}
\textsuperscript{(24)} tare da \VUgras{jakinta} \textsuperscript{(23)}.
Tana dawowa daga gari, inda ta kai madara mai kyau ga yara.

%%K t03-tit-5 sub
\VUsaut{4.0mm}
\VUpk{10.2pt}{40}{Gonar ƴaƴan itace.}
\VUsaut{3.4mm}

%%K t03-c1-09-1 p
9. --- Malam Jaki ya gamsu ƙwarai da wannan tafiya ta safe. Yana
shaƙar iskar da ƙamshin \VUgras{furanni} \textsuperscript{(20)} ya
ɗebi zaƙi, furannin da suka rufe \VUgras{itatuwan ƴaƴan itace}
\textsuperscript{(18)} na \VUgras{gonar} maƙwabta. Waɗannan
\VUgras{itatuwan fis} \textsuperscript{(18)} da \VUgras{itatuwan
abirkot} \textsuperscript{(19)} duk ruwan hoda da fari suke, kuma suna
ba da tsammanin girbi mai kyau.

%%K t03-c1-10-1 p
10. --- A halin yanzu, ƙananan \VUgras{ƙudan zuma} \textsuperscript{(22)}
na \VUgras{amyar zuma} \textsuperscript{(21)} suna zuwa can su tara
\VUgras{zumarsu} mai zaƙi suna ƙara.

%%K t03-c1-11-1 p
11. --- Amma me ke faruwa ? Ga shi yaran ƙauyen sun yi gudu suka
tsoratar da \VUgras{manyan agwagwa} \textsuperscript{(25)} suka sa su
gudu. Fitowar jama’a ce daga coci bayan \VUgras{baftisman} ɗan
magatakarda.

%%K t03-c1-12-1 p
12. --- Ƙaramin Iakobus, cikin \VUgras{shimfiɗarsa}
\textsuperscript{(26)}, ya farka saboda kaɗawar ƙararrawa mai farin
ciki, ƙanwarsa Magdalena kuwa, ita ma tana son ganin jama’ar
baftisma, sai ta roƙi mahaifiyarta ta zo ta yi mata kitso ta sa mata
tufafi a bakin ƙofar gida.

%%K t03-c1-13-1 p
13. --- Mahaifiyar ta yarda. Ta sa \VUgras{ɗankwalinta}
\textsuperscript{(28)}, da \VUgras{rigar jikinta} \textsuperscript{(29)}
da \VUgras{afaronta} \textsuperscript{(30)}, sai ta zo ta zauna a kan
ƙaramar kujera gaban ƙofa. Magdalena kuwa ba ta da komai sai
\VUgras{rigar ƙasanta} \textsuperscript{(31)} da \VUgras{matsattsiyar
rigarta} \textsuperscript{(32)}.

%%K t03-c1-14-1 p
14. --- Kai, farin ciki gare ta ! Ta manta da furanninta da take so,
wato \VUgras{furen kanumfari} \textsuperscript{(34)}, wanda
\VUgras{malam buɗe littafi} \textsuperscript{(35)} yake sauka a kansa,
da \VUgras{tulifinta} \textsuperscript{(36)}, da \VUgras{furannin
Mayunta} \textsuperscript{(37)} cikin \VUgras{tukwane}
\textsuperscript{(38)} a kan \VUgras{bango} \textsuperscript{(33)}, da
\VUgras{wardinta} \textsuperscript{(39)} da suka yi shinge mai kyau
ƙwarai. Tana kallon tufafi masu tsada na matan da suke cikin jama’a.

%%K t03-c1-15-1 p
15. --- Yaran ƙauyen ba su damu sosai da tufafi ba. Suna gaggawa suna
tunkuɗar juna domin su sami \VUgras{alewa} \textsuperscript{(55)} ko
\VUgras{tsabar kuɗi} \textsuperscript{(52)}. Ɗaya daga cikinsu yana
kuka \textsuperscript{(40)}.

%%K t03-c1-16-1 p
16. --- Wani ya taka ƙafar \VUgras{kare} \textsuperscript{(42)}, sai
kare ya gudu yana kuka ya tsoratar da \VUgras{kaza}
\textsuperscript{(43)} da \VUgras{ƴaƴanta} \textsuperscript{(44)}.

%%K t03-tit-6 sub
\VUsaut{5.0mm}
\VUpk{10.2pt}{40}{Jerin gwanon baftisma.}
\VUsaut{4.0mm}

%%K t03-c1-17-1 p
17. --- A gaban jama’a akwai \VUgras{mai reno} \textsuperscript{(45)}.
Tana da \VUgras{kyakkyawar hula} \textsuperscript{(46)} mai dogayen
ɗauri. Tana ɗauke da \VUgras{jaririn da aka haifa}
\textsuperscript{(47)} sanye da \VUgras{lesi} \textsuperscript{(48)} gaba
ɗaya.

%%K t03-c1-18-1 p
18. --- Bayan mai reno sai \VUgras{uban baftisma}
\textsuperscript{(49)} da \VUgras{uwar baftisma} \textsuperscript{(53)}.
Su ne suke rarraba alewa da tsabar kuɗi. Uban baftisma yana da
\VUgras{safar hannu} \textsuperscript{(51)} da \VUgras{dogon hula}
\textsuperscript{(50)}. Uwar baftisma tana riƙe da \VUgras{gungun
furanni} \textsuperscript{(54)} mai kyau a hannu.

%%K t03-c1-19-1 p
19. --- Bayansu, muna ganin \VUgras{baffan} \textsuperscript{(56)} da
\VUgras{goggon} \textsuperscript{(58)} yaron. Goggo tana da
\VUgras{hula} \textsuperscript{(59)} mai kwarjini, baffa kuwa yana da
\VUgras{bakar doguwar riga} \textsuperscript{(57)} da farin waskoti. Ga
kuma \VUgras{ɗan baffa} \textsuperscript{(60)} da \VUgras{ƴar baffa}
\textsuperscript{(61)}. Ita kuwa tana riƙe da hannun \VUgras{ƙanwar}
\textsuperscript{(62)} jaririn, wato ƙaramar Berta.

%%K t03-c1-20-1 p
20. --- \VUgras{Ɗan uwan} \textsuperscript{(63)} Berta ɗayan yana
tafiya a baya. Yana riƙe da hannun wata \VUgras{ƴar uwa}
\textsuperscript{(64)}. \VUgras{Abokan} \textsuperscript{(65)} iyalin
suna zuwa a bayansu.

%%K t03-tit-7 sub
\VUsaut{4.6mm}
\VUpk{10.2pt}{40}{Masu kallo.}
\VUsaut{3.4mm}

%%K t03-c1-21-1 p
21. --- Kusa da jama’a, ga \VUgras{maroƙi} \textsuperscript{(66)} yana
jingina a kan \VUgras{sandar tallafinsa} \textsuperscript{(67)}. Yana
neman sadaka.

%%K t03-c1-22-1 p
22. --- A ɗaya gefen akwai ƙaramin yaro \VUgras{mai ladabi}
\textsuperscript{(68)} ƙwarai. Yana gaisuwa yana yi wa waɗanda ya sani
« \VUgras{barka da safiya} ».

%%K t03-c1-23-1 p
23. --- Mutanen ƙauyen sun fito daga gidajensu domin su ga jerin
gwanon baftisma. Muna ganin \VUgras{ɗan ƙauye} \textsuperscript{(69)}
sanye da \VUgras{hular audugarsa,} a gefensa kuma \VUgras{ƴar ƙauye}
\textsuperscript{(70)}. Sai \VUgras{tsohuwa} \textsuperscript{(71)}
jingine a kan \VUgras{sandarta} \textsuperscript{(72)}. A ƙarshe
\VUgras{mai niƙa} \textsuperscript{(73)} da \VUgras{babbar hular
gashinsa} \textsuperscript{(74)} da wata budurwar ƙauye sanye da
\VUgras{takalman katako} \textsuperscript{(75)}, wadda take koya wa
ɗanta tafiya.

%%K t03-c1-24-1 p
24. --- Wannan fage yana da daɗi ƙwarai.

\VUsaut{4.0mm}
%%K t03-tit-8 sub
\VUcentre{12.0pt}{0}{\textit{Fage na Biyu.}}
\VUsaut{3.0mm}

%%K t03-tit-9 sub
\VUpk{10.2pt}{40}{Bikin jama’a.}
\VUsaut{2.4mm}

%%K t03-tit-10 sub
\VUpk{10.2pt}{40}{Wasannin ruwa da gasa.}
\VUsaut{3.0mm}

%%K t03-c2-01-1 p
1. --- \VUgras{Rani} ya zo. Zafin rana mai tsanani na watan
\VUgras{Yuni} (Yuli ko \VUgras{Agusta}) yana ƙarfafa samari su yi iyo
su kuma tuƙa jirgin ruwa.

%%K t03-c2-02-1 p
2. --- A gaɓar dama ta kogi, masu tuƙa jirgin ruwa suna tuɓe tufafi
domin su shiga wasannin ruwa da gasa.

%%K t03-c2-03-1 p
3. --- Ɗaya daga cikinsu yana cire \VUgras{takalmansa}
\textsuperscript{(1)} ; yana kwance su. Abokansa biyu sun riga sun yi
shiri. Yanzu ba su da komai sai \VUgras{rigar saƙa}
\textsuperscript{(2)} da \VUgras{wandon iyo} \textsuperscript{(3)}. Suna
hira suna jiran abokansu. Suna magana a kan hutu da bikin da yake
gudana a ɗaya gaɓar.

%%K t03-c2-04-1 p
4. --- A ƙasa, kusa da su, ga tufafin abokansu a kwance :
\VUgras{takalman sama} \textsuperscript{(4)} guda biyu,
\VUgras{takalma} \textsuperscript{(5)}, \VUgras{safa}
\textsuperscript{(6)}, \VUgras{hular gashi} \textsuperscript{(7)} da
\VUgras{hular kaba} \textsuperscript{(8)} da \VUgras{taye}
\textsuperscript{(9)}.

%%K t03-c2-05-1 p
5. --- Ɗaya daga cikin samarin ya naɗe \VUgras{kayansa}
\textsuperscript{(10)} da kula. Ya ajiye su a kan tawul, wanda
maƙwabcinsa zai naɗe kayansa biyu a cikinsa ba da daɗewa ba.

%%K t03-c2-06-1 p
6. --- Yaro na shida ya yi jinkiri. Har yanzu yana da
\VUgras{kasketinsa} \textsuperscript{(11)} a kai. Bai tuɓe
\VUgras{taguwarsa} \textsuperscript{(12)} ba, ko \VUgras{kwalar
rigarsa} \textsuperscript{(13)}, ko \VUgras{hannun rigarsa}
\textsuperscript{(14)}. Yana riƙe da \VUgras{waskotinsa}
\textsuperscript{(17)} a hannu, har yanzu kuma yana da
\VUgras{wandonsa} \textsuperscript{(16)} da \VUgras{ɗamararsa}
\textsuperscript{(15)}. Lalle ba zai yi shiri ba don gasar da take
gabatowa.

%%K t03-c2-07-1 p
7. --- Kusa da samarin, akwai ƙaramar hanya wadda a gefenta akwai
\VUgras{ciyawar ƙaya} \textsuperscript{(18)} mai yawa, tana kaiwa
bakin kogi, inda malam Iakobus yake shirya, cikin \VUgras{kyauro}
\textsuperscript{(19)}, \VUgras{wasan wuta} \textsuperscript{(20)} da za
a harba da yamma.

%%K t03-tit-11 sub
\VUsaut{4.4mm}
\VUpk{10.2pt}{40}{Gasar.}
\VUsaut{3.4mm}

%%K t03-c2-08-1 p
8. --- A kan kogi, waɗansu mata, suna kare kansu da laimarsu, suna
yawo cikin \VUgras{jirgin ruwa} \textsuperscript{(21)}. Suna bin gasar
da take da ban sha’awa ƙwarai. A cikin kowane \VUgras{ƙaramin jirgi}
\textsuperscript{(22)}, \VUgras{masu tuƙi} \textsuperscript{(24)} huɗu
suna tuƙi da dukan ƙarfinsu, alhali \VUgras{mai riƙe da linzami}
\textsuperscript{(26)} yana riƙe da \VUgras{linzamin}
\textsuperscript{(25)} yana kuma jagorantar jirgin nan mai rauni.
\VUgras{Sandunan tuƙi} \textsuperscript{(23)} suna bugun ruwa cikin
kari, jiragen nan marasa nauyi kuwa suna tafiya da sauri mai ban
mamaki.

%%K t03-c2-09-1 p
9. --- Waɗansu samari suna gwagwarmaya, kusa da ginshiƙin gada, a
kan ladar \VUgras{sandar mai} \textsuperscript{(28)}. Ɗaya daga
cikinsu yana takawa a hankali a kan sandar nan mai laushi da santsi
domin ya kai ga ƙaramar \VUgras{tuta} \textsuperscript{(29)} da aka
ɗaura a ƙarshenta. \VUgras{Abokin takararsa} \textsuperscript{(30)}
mara sa’a ya faɗa cikin ruwa. Yana iyo domin ya fito ƙasa.

%%K t03-c2-10-1 p
10. --- Samari manya suna \VUgras{gasar tunkuɗa a ruwa}
\textsuperscript{(32)} kusa da \VUgras{tashar ruwa}
\textsuperscript{(33)}. Waɗannan wasanni duka ana kallonsu da
\VUgras{jama’a} \textsuperscript{(31)} mai yawa, jingine a kan
\VUgras{shingen} \VUgras{gada} \textsuperscript{(27)} kuma taruwa a kan
gaɓa. Sai dai waɗansu masu kallo sukan juya domin su bi wasan
\VUgras{sandar lada} \textsuperscript{(45)} ko domin su yaba wa
\VUgras{tawagar magajin gari} wadda take zuwa domin \VUgras{rabon}
kyaututtuka.

%%K t03-c2-11-1 p
11. --- Da farko, ga waɗansu \VUgras{yaran gari}
\textsuperscript{(35)} a gaban \VUgras{makaɗa} \textsuperscript{(36)} ;
sai \VUgras{magajin gari} \textsuperscript{(37)} ya biyo, da
mataimakansa da \VUgras{majalisar} ƙaramin garin. A ƙarshe sai
\VUgras{masu kashe gobara} \textsuperscript{(38)}.

%%K t03-tit-12 sub
\VUsaut{5.0mm}
\VUpk{10.2pt}{40}{Bikin gari.}
\VUsaut{3.6mm}

%%K t03-c2-12-1 p
12. --- Mutane da yawa suna kan \VUgras{filin biki}
\textsuperscript{(39)}. Ana zuwa a ga \VUgras{rumfuna} iri iri :
\VUgras{dawakin katako} \textsuperscript{(40)}, da \VUgras{sirkas}
\textsuperscript{(41)} wanda wasansa ya riga ya fara ; da \VUgras{wurin
rawa} \textsuperscript{(42)}, inda samari da ƴan mata na garin suke
jin daɗin \VUgras{rawa.}

%%K t03-c2-13-1 p
13. --- A can baya, muna ganin \VUgras{filin fama}
\textsuperscript{(44)} inda \VUgras{mai fama da bijimi} na Sifaniya
mai jaruntaka yake kokawa da \VUgras{bijimi} \textsuperscript{(45)} mai
fushi ; sai \VUgras{wurin zamiya} \textsuperscript{(49)}, da
\VUgras{wurin harbi} \textsuperscript{(46)} inda ake horar da kai a kan
amfani da \VUgras{bindigogi} da harbi a kan \VUgras{abin harbi.}

%%K t03-c2-14-1 p
14. --- A kusa, ga \VUgras{gidan wasan kwaikwayo}
\textsuperscript{(47)} inda ake yin \VUgras{wasan ban dariya} da
\VUgras{wasan ban tausayi.} Kusa ƙwarai, akwai rumfar
\VUgras{katon mutum} \textsuperscript{(48)} da ta \VUgras{gajeren
mutum.} \VUgras{Tsawon} na farkon mita biyu da santimita goma sha
biyar ne ; na biyun kuwa bai fi yaro girma ba.

%%K t03-c2-15-1 p
15. --- A ƙarshe, ga babban \VUgras{gidan namun daji}
\textsuperscript{(50)} tare da \VUgras{namun dajinsa,} da
\VUgras{damisarsa} \textsuperscript{(51)} mai \VUgras{farata} masu
ƙarfi, da \VUgras{zakinsa} \textsuperscript{(52)} mai \VUgras{geza} mai
kauri, da \VUgras{raƙuman dawansa} \textsuperscript{(53)} biyu da
\VUgras{giwarsa} \textsuperscript{(54)} wadda take amfani da
\VUgras{hancinta} cikin gwaninta.

%%K t03-c2-16-1 p
16. --- \VUgras{Mai horar da namun daji} \textsuperscript{(55)} wanda
yake koya wa waɗannan dabbobi yana tsaye a bakin ƙofar rumfar.

\VUsaut{6.0mm}
\VUfilet{22mm}
